%!TEX root = ../main_thesis.tex

\chapter{Risoluzione dei modelli}
\label{cap:solver}

Giunti in fondo alla pipeline di compilazione, il modello è pronto per essere
risolto. Questo capitolo descrive le tecniche di risoluzione adottate da ROOC:
il simplesso a due fasi, implementato direttamente nel progetto a fini
didattici; il solver \texttt{microlp}, contributo originale di questo lavoro, che
estende il simplesso al caso intero; e i backend esterni, selezionati
automaticamente in base alla natura del modello.

\section{Concetti teorici}
\label{sec:solver-teoria}

\paragraph{Il metodo del simplesso.}
Come accennato nel Capitolo~\ref{cap:background}, la soluzione ottima di un
problema di programmazione lineare, se esiste, si trova in un vertice della
regione ammissibile. Il \emph{metodo del simplesso} sfrutta questa proprietà
spostandosi di vertice in vertice lungo gli spigoli del poliedro, scegliendo a
ogni passo una direzione che migliora il valore della funzione obiettivo. Su un
modello in forma standard, l'algoritmo lavora su un \emph{tableau} e procede per
\emph{pivot}: a ogni iterazione seleziona una variabile entrante (quella che
promette il miglioramento maggiore, individuata dai \emph{costi ridotti}) e una
variabile uscente (determinata dal \emph{test del rapporto minimo}, che preserva
l'ammissibilità), fino a quando non è più possibile migliorare e si è raggiunto
l'ottimo.

\paragraph{Il metodo a due fasi.}
Il simplesso richiede una base ammissibile di partenza, che non sempre è
immediatamente disponibile. Il \emph{metodo a due fasi} risolve questo problema
introducendo, in una prima fase, delle \emph{variabili artificiali} e
minimizzando la loro somma: se il minimo raggiunto è zero, si è ottenuta una
base ammissibile per il problema originale; la seconda fase ottimizza quindi la
funzione obiettivo vera e propria a partire da tale base.

\section{Applicazione in ROOC: il simplesso a due fasi}
\label{sec:simplex}

ROOC include un'implementazione propria del metodo del simplesso, costruita sopra
la forma standard prodotta nel Capitolo~\ref{cap:linearizzazione}. La struttura
centrale è il \texttt{Tableau}, che mantiene i coefficienti dell'obiettivo, la
matrice e il vettore dei termini noti dei vincoli, l'insieme delle variabili in
base e il valore corrente. A partire dalla forma standard il solver costruisce il
tableau, cercando una base ammissibile immediata e, quando questa non esiste,
ricorrendo al metodo a due fasi descritto in precedenza; esegue quindi i vari
passaggi del simplesso, iterando i pivot, fino a raggiungere l'ottimo, rilevare
l'illimitatezza del problema o toccare un limite massimo di iterazioni.

\paragraph{Risoluzione passo-passo.}
Oltre alla risoluzione diretta, il tableau offre una modalità \emph{passo-passo}
che, a ogni iterazione, ne registra lo stato prima del pivot insieme alle
variabili entrante e uscente: è questa la funzionalità che alimenta la
visualizzazione didattica della piattaforma (Capitolo~\ref{cap:piattaforma}).
Trattandosi di un'implementazione orientata alla chiarezza più che alle
prestazioni, questo simplesso è impiegato soprattutto a scopo illustrativo, mentre
la risoluzione efficiente è affidata ai solver descritti nelle sezioni seguenti.

\section{\texttt{microlp}: fork ed estensione a MILP}
\label{sec:microlp}

\texttt{microlp} è un solver pubblicato come \emph{crate} Rust indipendente. Nasce
come \emph{fork} di \texttt{minilp}, un crate ormai archiviato che implementava il
solo metodo del simplesso, limitato a variabili continue. Nell'ambito di questo
lavoro il crate è stato ripreso e mantenuto, ed esteso con il supporto a variabili
\emph{intere} e \emph{binarie}: da risolutore di sola programmazione lineare è
diventato così un risolutore di programmazione lineare intera mista (MILP). Questa
estensione costituisce uno dei contributi originali della tesi. Il crate ha
peraltro una diffusione che va oltre questo progetto e sta raccogliendo
l'interesse della comunità. È sviluppato come progetto \emph{open source} su
GitHub\footnote{\url{https://github.com/Specy/microlp}} ed è pubblicato su
\emph{crates.io},\footnote{\url{https://crates.io/crates/microlp}} dove, al
momento della scrittura, ha superato il milione di download; è inoltre utilizzato
da \texttt{good\_lp}, tra le più note librerie di programmazione lineare per Rust.
Attorno al progetto si è formata una comunità che contribuisce attivamente,
segnalando problemi e proponendo miglioramenti tramite \emph{issue} e
\emph{pull request}.

\paragraph{Estensione al caso intero: branch-and-bound.}
Il simplesso risolve problemi a variabili continue, ma non garantisce soluzioni
intere. La tecnica con cui \texttt{microlp} affronta l'integralità è il
\emph{branch-and-bound}: si risolve dapprima il \emph{rilassamento continuo} del
problema, ignorando i vincoli di integralità; se una variabile che dovrebbe essere
intera assume un valore frazionario, si \emph{ramifica} il problema in due
sottoproblemi, imponendo nell'uno che la variabile sia minore o uguale alla parte
intera inferiore del valore e nell'altro che sia maggiore o uguale a quella
superiore. Il valore del rilassamento fornisce inoltre una \emph{limitazione}
(\emph{bound}) che consente di potare i rami che non possono contenere soluzioni
migliori di quella già trovata.

All'interno di ROOC, \texttt{microlp} svolge un duplice ruolo: risolve i modelli
puramente continui, come alternativa al backend esterno descritto più avanti, e
soprattutto è il motore impiegato per tutti i modelli che contengono variabili
intere o miste. La risoluzione avviene costruendo un problema nel formato del
solver: a ogni variabile del modello viene associata una variabile del solver,
con il proprio dominio (reale, intero o binario) e i propri estremi; vengono
quindi aggiunti i vincoli, tradotti nei rispettivi operatori di confronto, e
infine si richiede la soluzione, da cui si estraggono i valori delle variabili e
il valore ottimo.

\section{Backend esterni e \emph{auto-solver}}
\label{sec:backend-esterni}

Oltre al solver proprio, ROOC integra alcuni backend esterni, ciascuno
specializzato per una particolare classe di problemi. Per i modelli a sole
variabili continue è disponibile il solver \emph{Clarabel}, utilizzato tramite la
libreria \texttt{good\_lp}. I modelli con variabili intere o binarie vengono invece
risolti tramite il solver integrato \texttt{microlp}.

\paragraph{Selezione automatica del solver.}
La presenza di più solver, ciascuno adatto a un diverso tipo di modello, è
gestita in modo trasparente da un \emph{auto-solver}, che ispeziona il dominio
delle variabili e instrada il problema verso il risolutore appropriato.
